% Part: counterfactuals
% Chapter: minimal-change-semantics
% Section: antecedent-strengthening

\documentclass[../../../include/open-logic-section]{subfiles}

\begin{document}

\olfileid{cnt}{min}{agg}

\olsection{पूर्वांग बळकट करणे}

‘‘पूर्वांग बळकट करणे’’ म्हणजे
$!A \lif !C \Entails
(!A \land !B) \lif !C$
हे अनुमान. वास्तविक
अभिव्यंजनासाठी ते वैध,
पण प्रतिवास्तविक
विधानांसाठी अवैध आहे.
मी ही आगकाडी ओढली असती,
तर ती पेटली असती,
हे सत्य आहे असे समजा.
(म्हणजे आगकाडीत किंवा
आगपेटीच्या घासण्याच्या
पृष्ठभागात काही दोष नाही,
मी आगकाडी मोडणार नाही,
इत्यादी.) पण ही
आगकाडी मी बाह्य
अवकाशात ओढली असती,
तर ती पेटली असती,
हे सत्य नाही. म्हणून
पुढील अनुमान अवैध आहे:
\begin{quote}
  आगकाडी ओढली असती,
  तर ती पेटली असती.

  म्हणून आगकाडी बाह्य
  अवकाशात ओढली असती,
  तर ती पेटली असती.
\end{quote}

लुईस--स्टालनाकर यांच्या
सशर्त विधानांच्या
विश्लेषणातून हे
समजते: मी आगकाडी
बाह्य अवकाशात
ओढतो असे सर्वाधिक
जवळचे जग, मी
आगकाडी ओढतो
अशा सर्वाधिक
जवळच्या जगापेक्षा
वास्तविक जगापासून
खूप दूर आहे.
म्हणून दुसऱ्या
जगात आगकाडी
पेटते हे सत्य
असले, तरी
पहिल्या जगात
ती पेटत नाही.
असेच असायला हवे.

\begin{ex}
  गोलक चिन्हार्थमीमांसेत हे अनुमान
  अवैध ठरते; म्हणजे
  $p \cif r \Entails/
  (p \land q) \cif r$.
  $\mModel{M} = \tuple{W, O, V}$
  हे प्रतिमान पाहा: येथे
  $W = \{w, w_1, w_2\}$,
  $O_w = \{\{w\}, \{w,
  w_1\}, \{w, w_1, w_2\}\}$,
  $V(p) = \{w_1, w_2\}$,
  $V(q) = \{w_2\}$, आणि
  $V(r) = \{w_1\}$.
  इतर जगांसाठी $O_{w_1}=\{\{w_1\}\}$ आणि
  $O_{w_2}=\{\{w_2\}\}$ घ्या; त्यामुळे $O$ हे
  सर्व $W$ वर परिभाषित आहे.
  $p$ समाविष्ट करणारा
  $S = \{w, w_1\}$
  हा गोलक आहे आणि
  त्यातील सर्व जगांत
  $p \lif r$ सत्य आहे;
  म्हणून
  $\mSat{M}{p \cif r}[w]$.
  तसेच $(p \land q)$
  समाविष्ट करणारा
  $S' = \{w, w_1, w_2\}$
  हा गोलक आहे; पण
  $\mSat/{M}{(p \land q) \lif r}[w_2]$.
  $S'$ हा या $O_w$ मधील एकमेव $(p\land q)$-जग
  धारण करणारा गोलक आहे आणि त्यात अभिव्यंजन असत्य आहे.
  म्हणून कोणताही साक्षी-गोलक नाही; रिक्तपणे सत्य होण्याची
  अटही लागू होत नाही. त्यामुळे
  $\mSat/{M}{(p \land q) \cif r}[w]$
  (\olref{fig:antecedent-strengthening} पाहा).

\begin{figure}
\begin{center}
\begin{tikzpicture}[layerwidth=1.5,scale=.6]\tiny
  \spheresystem{3}
  \propositionintersect{-10}{3}{30}{5}
  \begin{scope}
    \clip plot[smooth,tension=1] coordinates
          {(0,4.5) (30:.95) (4.5,-3.5)} -- (4.5,4.5) ;
    \spherefill{2}
  \end{scope}
  \draw[proposition] plot[dashed,smooth,tension=1] coordinates
          {(0,4.5) (30:.95) (4.5,-3.5)} ;
  \proposition{30}{2}{30}{5}
  \path (0,0) node[world] {$w$};
  \spherepos{30}{2}{node[world] {$w_1$}}
  \spherepos{-10}{3}{node[world] {$w_2$}}
  \spherepos{-10}{4}{node {$q$}}
  \spherepos{30}{4}{node {$r$}}
  \draw (1,4.5) node[below] {$p$};
\end{tikzpicture}
\caption{पूर्वांग बळकट करणे या नियमाचे प्रतिउदाहरण}
\ollabel{fig:antecedent-strengthening}
\end{center}
\end{figure}
\end{ex}

\end{document}
